\section{Future Work}
\label{sec:future}

This section maps research extensions to the conservative trim algorithm,
prioritised by estimated reward-per-effort. All directions build on the
frequency-weighted disclosure framework introduced here.

\paragraph{Learned compression with instruction-document benchmarks.}

The literature gap identified in \cref{sec:related} is concrete and addressable:
build a benchmark suite of instruction documents (from \blackrim, AutoGen,
CrewAI, and other multi-agent frameworks) with fidelity labels, then compare
conservative trim (this work) against LLMLingua-style learned compression.
The benchmark would measure token savings, instruction-fidelity preservation,
and latency (to capture the cost of dynamic disclosure). Such a benchmark
would settle whether learned compression is applicable to procedural documents
or whether conservative trim is the right tool for the job.

\paragraph{Dynamic loading via tool calls.}

Externalised sections could be retrieved on-demand rather than externalised
away entirely. Introduce a \texttt{read\_instruction(section\_id)} tool that
agents invoke when they need a backgrounder on (say) worktree isolation.
The trade-off is clear: fewer tokens in the static prefix (20--30\,\% additional
savings beyond this work) but variable latency per spawn (one extra RPC
round-trip when a section is first touched). Measure this on real agent
workflows to determine the latency/token cost curve and whether
dynamic loading is a net win for different agent types.

\paragraph{Cross-session prefix sharing in multi-instance deployments.}

When multiple instances of \blackrim run in parallel (or across days), the
\claudemd prefix is reloaded and re-cached for each instance. Anthropic's
documentation suggests prompt-cache can be shared across API keys under
certain conditions. Explore whether a multi-instance deployment (e.g.,
staging and prod, or two independent teams using Blackrim) can share
the same cached prefix on-disk or in Anthropic's cache infrastructure,
reducing first-spawn cost to near-zero per instance.

\paragraph{Per-section frequency estimation from session logs.}

The current work hand-classifies sections as \emph{frequent}, \emph{occasional},
or \emph{rare} using domain knowledge and the author's intuition. A more
principled approach: instrument agent spawns to track which sections were
semantically used (not just loaded) per spawn, then build a lived frequency
histogram. This is non-trivial (requires shallow semantic instrumentation
of agent reasoning) but would remove author bias from the trim decision.

\paragraph{Multi-document residency: CLAUDE.md + AGENTS.md + RTK.md.}

Extend the model from a single instruction document to the full suite:
CLAUDE.md (framework rules), AGENTS.md (citizen/worker definitions), and
any runtime knowledge base (RTK.md). This exposes new optimisation surface:
inter-document cross-references complicate the externalisation decision
(a section of AGENTS.md might be referenced from CLAUDE.md, making simple
externalisation risky). Develop a multi-document dependency graph and
apply the conservative trim algorithm across the whole set.

\paragraph{Application beyond Blackrim: support agents, retrieval, research.}

This paper's methodology was motivated by and validated on a multi-agent
coding framework. Does the conservative trim model transfer to other instruction-heavy
domains? Test the algorithm on (1) customer-support agents (knowledge bases
for FAQ, escalation rules), (2) retrieval-augmented generation systems (search
result ranking, reranking instructions), and (3) research-automation agents
(literature review, experimental design). Each domain has different duplication
patterns and fidelity constraints, so the three would jointly establish scope.
